\documentclass[12pt]{article}
\usepackage[utf8]{inputenc}
\usepackage{amsmath}
\usepackage{amssymb}
\usepackage{geometry}
\usepackage{listings}
\usepackage{xcolor}
\usepackage{hyperref}

\geometry{a4paper, margin=1in}

\definecolor{codegreen}{rgb}{0,0.6,0}
\definecolor{codegray}{rgb}{0.5,0.5,0.5}
\definecolor{codepurple}{rgb}{0.58,0,0.82}
\definecolor{backcolour}{rgb}{0.95,0.95,0.92}

\lstset{
    backgroundcolor=\color{backcolour},   
    commentstyle=\color{codegreen},
    keywordstyle=\color{magenta},
    stringstyle=\color{codepurple},
    basicstyle=	tfamily\small,
    breaklines=true,                 
    numbers=left,                    
    numbersep=5pt,                  
    tabsize=2
}

	itle{Documentation: AO* Search Algorithm for AND-OR Graphs}
\author{AI Lab - Assignment 3}
\date{	oday}

\begin{document}

\maketitle

\section{Introduction}
The AO* algorithm is used to find the optimal solution path in an AND-OR graph. Unlike standard state-space search (OR-graphs only), AND-OR graphs represent problems that can be solved by either one of several alternatives (OR) or by solving a set of sub-problems (AND).

\section{Mathematical Concepts}
In AO*, we don't use a $g(n)$ cost. Instead, we rely on updated heuristic values $h(n)$ that propagate upwards from leaf nodes to the root.

\subsection{Cost Calculation}
For a node $v$, let $S_1, S_2, \dots, S_k$ be the sets of successor nodes (branches). Each branch can be an AND branch (multiple nodes) or an OR branch (single node). The cost of a branch $S_i$ is calculated as:
\begin{equation}
    Cost(S_i) = \sum_{n \in S_i} [c(v, n) + h(n)]
\end{equation}
Where $c(v, n)$ is the cost of the edge (typically 1).

The heuristic value $h(v)$ is updated using the 	extbf{min-cost} branch:
\begin{equation}
    h(v) = \min_{i} \{ Cost(S_i) \}
\end{equation}

\section{Algorithm Logic}
The AO* algorithm in 	exttt{search\_graph.py} follows these steps:
\begin{enumerate}
    \item 	extbf{Initial Heuristic}: Start with initial heuristic estimates for all nodes.
    \item 	extbf{Expansion}: Expand the current node and identify its AND/OR successors.
    \item 	extbf{Bottom-up Update}: Once a leaf is reached or a node is expanded, recalculate the $h$ values of its parents based on the formula above.
    \item 	extbf{Solution Graph}: Maintain a 	exttt{solution\_graph} dictionary that tracks which branches currently yield the minimum cost.
\end{enumerate}

\section{Implementation Details}
\begin{itemize}
    \item 	exttt{compute\_minimum\_cost\_child\_nodes(v)}: This function iterates through all possible hyper-edges (branches) from $v$ and returns the sum of $(1 + h(child))$ for the cheapest branch.
    \item 	exttt{ao\_star(v, backtracking)}: Recursively processes nodes. When a node's heuristic is updated, it triggers updates in its ancestors (backtracking) to ensure global optimality.
\end{itemize}

\section{Graph Representation}
The graph is represented as a dictionary where keys are nodes and values are lists of lists. Each inner list represents a branch:
\begin{itemize}
    \item 	exttt{A: [[B, E], [D]]} means node A has an AND-branch to \{B, E\} and an OR-branch to \{D\}.
\end{itemize}

\end{document}
